\documentclass[12pt]{article}
\usepackage[T1]{fontenc}
\usepackage[utf8]{inputenc}
\usepackage{lmodern}
\usepackage{textcomp}
\usepackage{enumitem}
\usepackage{geometry}
\geometry{margin=1in}
\begin{document}
\begin{center}
Answer Key - Health Sciences - K-12 Level Assessment 4 \
\small (Answers are ordered exactly as in the question paper.)
\end{center}

\section*{Multiple Choice}
\begin{enumerate}[label=\arabic*.]
\item \textbf{Answer:} B
\\ \textit{Options:} A) 30 to 50\%; B) 40 to 60\%; C) 20 to 40\%; D) 10 to 30\%
\\ \textit{Note:} The correct answer of 40 to 60\% reflects the significant occurrence of alcoholic myopathy among individuals with alcohol use disorder, indicating that a substantial proportion of this population is affected by muscle weakness and damage linked to chronic alcohol consumption.
\item \textbf{Answer:} B
\\ \textit{Options:} A) 29; B) 32; C) 39; D) 43
\\ \textit{Note:} The correct answer is 32 because the body mass index (BMI) is calculated using the formula BMI = weight (kg) / (height ($m))^2$, which in this case is 99 kg / (1.75 $m)^2$, resulting in a BMI of approximately 32.16, rounded to 32.
\item \textbf{Answer:} D
\\ \textit{Options:} A) Goitre; B) Poor cognitive development; C) Poor bone growth; D) Increased risk of mortality
\\ \textit{Note:} Deficiency of vitamin A in children increases the risk of mortality because it compromises their immune system, making them more susceptible to infections and diseases that can lead to serious health complications.
\item \textbf{Answer:} C
\\ \textit{Options:} A) Chylomicrons; B) HDL; C) VLDL; D) LDL
\\ \textit{Note:} VLDL, or very low-density lipoprotein, is the primary lipoprotein secreted by the liver that transports dietary-derived lipids, particularly triglycerides, from the liver to various tissues in the body.
\item \textbf{Answer:} D
\\ \textit{Options:} A) Biotin; B) Niacin; C) Riboflavin; D) Vitamin C
\\ \textit{Note:} Vitamin C is essential for collagen synthesis because it acts as a cofactor for the enzymes that catalyze hydroxylation reactions, which are necessary for the stabilization and proper formation of collagen fibers.
\end{enumerate}

\section*{True / False}
\begin{enumerate}[label=\arabic*.]
\setcounter{enumi}{5}
\item \textbf{Answer:} False
\\ \textit{Options:} A) True; B) False
\\ \textit{Note:} The statement is false because different tissue proteins serve various functions and are regulated by distinct metabolic pathways, leading to varied rates of breakdown based on their specific roles in the body.
\item \textbf{Answer:} False
\\ \textit{Options:} A) True; B) False
\\ \textit{Note:} C$_{18}$ and C$_{22}$ polyunsaturated fatty acids are not the main precursors for eicosanoids; instead, it is the C$_{20}$ polyunsaturated fatty acid, arachidonic acid, that primarily serves as the precursor for this group of signaling molecules.
\item \textbf{Answer:} False
\\ \textit{Options:} A) True; B) False
\\ \textit{Note:} Crohn's disease can affect any part of the gastrointestinal tract, from the mouth to the anus, not just the colon.
\item \textbf{Answer:} False
\\ \textit{Options:} A) True; B) False
\\ \textit{Note:} The fear of becoming fat in patients with anorexia nervosa is not classified as a paranoid delusion because it is a distorted perception of body image and weight rather than a fixed false belief disconnected from reality.
\item \textbf{Answer:} True
\\ \textit{Options:} A) True; B) False
\\ \textit{Note:} Starch is a polysaccharide made up of long chains of glucose molecules, which oral bacteria lack the necessary enzymes to efficiently break down into simpler sugars for metabolism.
\end{enumerate}

\section*{Long Form Response}
\begin{enumerate}[label=\arabic*.]
\setcounter{enumi}{10}
\item \textbf{Answer:} Deficiencies in specific nutrients, particularly iron and copper, can lead to anemia, a condition where the body lacks enough healthy red blood cells to carry adequate oxygen. Iron is crucial for the production of hemoglobin, the protein in red blood cells that binds to oxygen; without enough iron, the body cannot produce sufficient hemoglobin, leading to fatigue and weakness. Copper plays a supportive role by helping to absorb iron and facilitating the formation of red blood cells. A lack of copper can hinder the body's ability to utilize iron effectively, further exacerbating anemia. Therefore, both iron and copper are essential for maintaining healthy blood function and preventing anemia.
\\ \textit{Note:} The answer correctly highlights that iron deficiency directly impairs hemoglobin production, essential for oxygen transport, while copper is necessary for iron absorption and red blood cell formation, demonstrating how both nutrients are vital in preventing anemia.
\item \textbf{Answer:} Nitrogen balance is an important concept in understanding how our bodies use protein. When a person is in a positive nitrogen balance, it means that the amount of nitrogen taken in from dietary proteins is greater than the amount excreted as waste. This state is crucial for growth, especially in children and adolescents, as it supports the development of new tissues, muscles, and overall health. Being in a positive nitrogen balance can also benefit recovering patients or athletes who need to rebuild their bodies after illness or intense training. Therefore, maintaining a positive nitrogen balance is essential for promoting growth and supporting overall health.
\\ \textit{Note:} correct because it accurately defines positive nitrogen balance as a state where nitrogen intake exceeds nitrogen excretion, highlighting its significance for growth, tissue development, and recovery in individuals.
\end{enumerate}

\vfill
\noindent\textit{End of Answer Key}
\end{document}